\section{Autoconvolution}%
\label{sec:autoconvolution}

Consider the autoconvolution operator $K$ in $L^2(0,1)$ defined for $x \in L^2(0,1)$ by
\[
  K(x)(t) = \int_0^t x(t-s)x(s)\,ds, \quad s \in (0,1)
\]
with linearization around $x_0 \in L^2(0,1)$ given by
\[
  K'(x_0)h(t) = 2\int_0^t x_0(t-s)h(s)\,ds, \quad t \in (0,1).
\]
\begin{enumerate}[label = (\alph*)]
  \item Take $x(s) = 2e^{-2s}$ and show that $y(t) = K(x)(t) = 4te^{-2t}$.
\end{enumerate}
\nopagebreak
\noindent\hrulefill

Taking \(x(s) = 2e^{-2s}\) and plugging it into the autoconvolution operator definition, we get
\[
  K(x)(t) = \int_{0}^{t} x(t - s) x(s) \odif{s} = \int_{0}^{t} 2e^{-2(t - s)} 2 e^{-2s} \odif{s} = 4e^{-2t}
  \int_{0}^{t} 1 \odif{s} = 4e^{-2t} \left[ t \right]_0^{t} = 4te^{-2t} = y(t)
\]
\nopagebreak
\noindent\hrulefill

\begin{enumerate}[resume, label = (\alph*)]
  \item Implement $K$ numerically using a quadrature rule as a mapping from $\mathbb{R}^{256}$ to $\mathbb{R}^{256}$.
    Use your code to find an approximation of $K(x)(t)$ and compare graphically to $y(t)$.
\end{enumerate}

\nopagebreak
\noindent\hrulefill

\begin{figure}[h]
  \centering
  \begin{subfigure}[t]{0.48\textwidth}
    \includegraphics[width=\textwidth]{fig_compare_quadrature_v_analytical_forward_autoconvolution.pdf}
    \caption{\label{fig:comparison-analytical-v-numerical-forward-autoconvolution}\
      Comparison of the quadrature rule numerical solution \(K(x)(t)\) to the analytical solution \(y(t)\).
    }
  \end{subfigure}\hfill
  \begin{subfigure}[t]{0.48\textwidth}
    \includegraphics[width=\textwidth]{fig_compare_quadrature_v_analytical_forward_autoconvolution_diff.pdf}
    \caption{\label{fig:compare-quadrature-v-analytical-forward-autoconvolution-difference}\
      The difference of the numerical and the analytical solutions for the autoconvolution operator.
    }
  \end{subfigure}
  \caption{\label{fig:comparison-forward-autoconvolution}\
    A graphical comparison of the numerical solution with the analytical solution of the autoconvolution operator.
  }
\end{figure}

\nopagebreak
\noindent\hrulefill

\begin{enumerate}[resume, label = (\alph*)]
  \item Add noise of relative size 5\% to $y$ to get $y_\delta$ and consider solving $K(x) = y_\delta$ by linearization:
    For the constant function $x_0(t) = 2 - 1.5t$, discretize the operator $K'(x_0)$ as a matrix $A_0 \in
      R^{256\times256}$.
    Solve using Tikhonov regularization the linearized and discretized inverse problem to find an approximate solution
      $x_1(t)$.
    Choose the regularization according to one of the parameter choice rules.
    Plot the solution and compare to $x(t)$. Calculate the error.
  \item Take another Newton step to find $x_2(t)$.
\end{enumerate}

\begin{figure}[h]
  \centering
  \begin{subfigure}[t]{0.48\textwidth}
    \includegraphics[width=\textwidth]{fig_first_newton_step.pdf}
    \caption{\label{fig:first-newton-step}\
      A visualization of the first newton step of the reconstruction. 
      The Tikhonov regularization parameter of \(\alpha = 0.05\) for the solution of \(h_0\) was chosen by empirically.
    }
  \end{subfigure}\hfill
  \begin{subfigure}[t]{0.48\textwidth}
    \includegraphics[width=\textwidth]{fig_newton_iterations_reg_lines.pdf}
    \caption{\label{fig:newton-iterations}\
      First \(4\) newton iterations of the linearized reconstruction.
      It is observable that the higher the iteration the greater the need for the regularization.
      This is sensible intuitively, as we are getting closer to the true solution and therefore also the relative error
        of the difference \(y_\delta - K(x_i)\) increases with the iteration number \(i\).
    }
  \end{subfigure}
  \caption{\label{fig:linearization-recontruction}\
    Newton iteration reconstruction.
  }
\end{figure}

\nopagebreak
\noindent\hrulefill

\begin{enumerate}[resume, label = (\alph*)]
  \item What happens if you start from the constant function $x_0(t) = -2 + 1.5t$ instead?
  \item Is the inverse problem $(K(x) = y)$ ill-posed; is it locally ill-posed?
\end{enumerate}

\nopagebreak
\noindent\hrulefill

Even if we increase the number of iterations, the reconstruction fails to converge as is visible in the
\cref{fig:newton-iterations-from-far-function}.
This can be regarded as a visual demonstration of local ill-posedness of the autoconvolution operator.

\begin{figure}[h]
  \centering
  \begin{subfigure}[t]{0.48\textwidth}
  \includegraphics[width=\textwidth]{fig_newton_iterations_reg_lines_alt.pdf}
  \caption{\label{fig:newton-iterations-from-far-function}\
    \(15\) Newton iterations of the inverse linearized autoconvolution reconstruction starting from a distant function from the
      true solution.
  }
  \end{subfigure}\hfill
  \begin{subfigure}[t]{0.48\textwidth}
  \includegraphics[width=\textwidth]{fig_newton_iterations_reg_lines_minus.pdf}
  \caption{\label{fig:newton-iterations-from-far-function-minus}\
    \(4\) Newton iterations of the inverse linearized autoconvolution reconstruction starting from a distant function
    from the desired solution converges to \(-x\).
  }
  \end{subfigure}
  \caption{\label{fig:newton-iterations-from-afar}\
    Newton iterations of the inverse linearized autoconvolution reconstruction starting from a distant function.
  }
\end{figure}

This also brings us to the subject of ill-posedness.
Due to the quadratic nature of the autoconvolution operator, if \(x\) is a solution to the inverse problem, then \(-x\)
  is also a solution with the same mapping.
With this in mind we can understand more, why an arbitrary number of iterations would not suffice for the newton steps
  from the initial linearization point function \(x_0 = -2 + 1.5t\) to converge to the desired solution \(x\).
  In fact it converges instead to the function \(-x\) as is visible from the
  \cref{fig:newton-iterations-from-far-function-minus}.
This is also a demonstration of the non-uniqueness of the convolution operator and this argumentation holds for any
  quadratic operator.
Due to the fact that the forward operator has a ``smoothing'' effect, we can intuitively expect local ill-posedness due
  to the inherent lack of continuousness with respect to small perturbations of the mapping \(y\).
In other words if we have \(y + \delta y \), we can expect the solutions error \(\delta x\) of the inverse problem,
  denoting \(\hat{x}(y + \delta y) = x + \delta x\), to be have a much larger norm than the that of \(\delta y\)
  (especially in any Sobolev norm).

\nopagebreak
\noindent\hrulefill

\section{Radial Calderón problem}%
\label{sec:radial-calderon-problem}

The radial FIT/Calderón problem linearized around $\gamma_0 = 1$ is defined through the operator $K$ as follows \[
  Kh(n) = -2\int_0^1 h(r)r^{2n-1}\,dr = \lambda_n, \tag{1}
\]
with $\lambda_n$ denoting eigenvalues of $\Lambda_\gamma - \Lambda_1$.
\begin{enumerate}[label = (\alph*)]
  \item In the file \texttt{DHW3.mat} on inside you find 10 precomputed (`measured') eigenvalues for a particular
      $\Lambda_\gamma - \Lambda_1$.
    Compute an estimate of $\gamma = 1 + h$ using both Tikhonov and TV regularization and plot it.
\end{enumerate}

\nopagebreak
\noindent\hrulefill

\begin{enumerate}[resume, label = (\alph*)]
  \item Knowing that $\gamma$ has the radial form
    \[
      \gamma(x) =
      \begin{cases}
        1 + \hat{q}, & |x| \leq R, \\
        1, & R < |x| \leq 1.
      \end{cases}
    \]
    give an estimate from the above of the contrast $\hat{q}$ and radius $R$.
\end{enumerate}
